\section{Introduction}

The $L^2$ metric on a moduli space of instantons records the size of
infinitesimal changes of a connection. Its behaviour under bubbling is
less direct: the curvature energy retains an atom at each bubble point,
whereas the tangent vectors are defined on smooth connections.
On a projective K3 surface this metric is hyperk\"ahler, but the sheaf
compactification and the Uhlenbeck compactification contract different
geometric data. The question is which compact space carries its limiting
metric, and how that metric restricts to the boundary.

We construct a local K\"ahler potential that is continuous at ideal
connections. It combines a zeta determinant on a curve avoiding the bubble
points with a logarithmically weighted curvature integral. Energy
quantization supplies an explicit correction for the defect cycle.
On K3 surfaces, this potential identifies the $L^2$ metric with the
canonical singular Ricci-flat metric in its determinant class. Completion,
volume and orbifold regularity then follow from the corresponding
algebraic and analytic theorems.

\subsection{The potential at an ideal connection}
Let $(X,\omega)$ be a smooth projective K\"ahler surface with
$[\omega]=c_1(H)$. We fix rank $r\ge2$, determinant and Chern classes.
For an ideal Hermitian--Yang--Mills connection $(F,Z)$, choose a smooth
curve $C\in|H^k|$ disjoint from $\Supp Z$, and a vector bundle $G$ on $C$
for which $\End F|_C\otimes G$ is acyclic. Such choices always exist.
If $g_C=k^{-1}\log|s_C|$ and $\rho=\rk G$, our potential is
\begin{equation}\label{eq:intro-potential}
 -\frac{\pi^2}{rk\rho}\log\det_\zeta\Delta_{\End F|_C\otimes G}
 -\frac12\int_X g_C|F_{A_F}^0|^2dV_X
 -4\pi^2\sum_x\operatorname{mult}_x(Z)g_C(x).
\end{equation}
Here $\ddc=(i/\pi)\partial\bar\partial$, the trace defining the metric
is the fundamental trace, and $\ddc g_C=k^{-1}[C]-\omega$.
The precise construction and coefficient calculation occupy
Sections~\ref{sec:potential}--\ref{sec:continuity}.

\begin{theorem}[Continuous potential]\label{thm:intro-potential}
The choices $C,G$ can be fixed in a neighbourhood of every ideal
connection. Formula~\eqref{eq:intro-potential} is continuous on that
neighbourhood in the Uhlenbeck topology, and its complex Hessian on the
smooth stable bundle locus is the $L^2$ K\"ahler form. On an obstructed
moduli space the same identity holds after pullback to each smooth
holomorphic family of stable bundles.
\end{theorem}

The curve serves two purposes. Its determinant gives the tangential
curvature term through the Quillen formula, and its logarithmic potential
transgresses that term to the surface. Since the curve avoids the bubbles,
the singular test function is integrated only where the curvature
converges smoothly. Away from the curve, weak convergence of measures is
sufficient. This separation is the reason a continuous potential can be
obtained without estimates for instanton gluing parameters.

\subsection{The K3 consequences}
Let now $X$ be a projective K3 surface with Ricci-flat K\"ahler form
$\omega\in c_1(H)$. Let $v=(r,c,s)$ be primitive, $r\ge2$, $v^2>0$,
and let $H$ be $v$-generic. Assume that the moduli space $\M$ of slope-stable
locally free sheaves is nonempty. Set $2N=v^2+2$ and
$V=\Vol(X,g_X)$. Let $\Mn$ be the normalization of the Uhlenbeck
compactification; it is canonically homeomorphic to that compactification.

\begin{theorem}[Canonical metric and boundary]\label{thm:intro-main}
The $L^2$ K\"ahler form extends to the canonical Ricci-flat current in
an ample determinant class on $\Mn$, with continuous local potentials.
Its metric completion is the Uhlenbeck compactification and
\[
 \Vol(\M,\gL)=\frac{(16\pi^4V)^N}{N!}.
\]
For rank two, along every stratum with stable residual bundle and $\ell$
distinct simple bubble points, the completed metric is a smooth
hyperk\"ahler orbifold metric. The stratum is totally geodesic, with
intrinsic metric the product of the lower-charge $L^2$ metric and
$\Sym^\ell(X,4\pi^2g_X)$ restricted to distinct points.
\end{theorem}

The current identification is proved in
Theorems~\ref{thm:global-current} and~\ref{thm:canonical}; the completion
and boundary statements are Theorems~\ref{thm:completion}
and~\ref{thm:orbifold}. They also give finite curvature energy and an
explicit lower bound for the diameter. A power of the Quillen metric
extends continuously to an ample line on $\Mn$. In contrast, the $L^2$
metric cannot extend as a continuous positive-definite metric to the
Gieseker moduli space when the Uhlenbeck morphism contracts a curve.

At a boundary point with stable residual bundle, we further identify the
affine tangent cone and its Reeb grading with a product of affine ADHM
instanton spaces. This determines its volume density by the corresponding
Nekrasov coefficient (Theorems~\ref{thm:cones} and~\ref{thm:density}).
Outside the rank-two distinct-point locus, identification of the entire
tangent metric with the flat quotient metric is stated conditionally in
Remark~\ref{rem:cone-uniqueness}.

\subsection{Relation to earlier work}
Feehan \cite[Theorem~1.1]{Fee95} proves finiteness of the $L^2$ volume
and diameter for irreducible ASD moduli spaces on closed, connected,
oriented, simply connected four-manifolds with a generic Riemannian
metric, assuming either $b^+=0$ with $G=\mathrm{SU}(2)$ or
$\mathrm{SO}(3)$, or $G=\mathrm{SO}(3)$ with $w_2(P)\ne0$.
Under the same hypotheses, \cite[Theorem~1.2]{Fee95} states that the
metric completion is homeomorphic to the Uhlenbeck compactification;
the proof of that statement is deferred there to a separate work.
These hypotheses and ours have different scopes. On a K3 surface the
$\mathrm{SO}(3)$ case with $w_2(P)\ne0$ is relevant, but the prescribed
Ricci-flat metric need not satisfy Feehan's generic-metric hypothesis.
Here the genericity assumption concerns the polarization for Gieseker
stability. For primitive positive-square Mukai vectors with nonempty
stable bundle locus, our argument applies in every rank $r\ge2$ to the
Ricci-flat representative itself. The continuous boundary potential
identifies its $L^2$ metric with the canonical singular Ricci-flat metric;
this gives the completion theorem in this setting, the exact volume,
and the rank-two orbifold boundary geometry stated above.

Groisser's survey \cite{Gro97} discusses the completion problem, and his analysis
of totally geodesic boundaries \cite{Gro98} establishes the scaled base
metric in charge-one settings and discusses higher-charge strata.
Our boundary formula retains the same factor
$4\pi^2$ on the base surface, now directly from the curvature-energy atom.

Fibre-integral expressions for the generalized Weil--Petersson form and
its determinant-line interpretation are established in
\cite{ST92,BS16}. Continuous potentials near locally free polystable
bundles are obtained in \cite{BSc20}.
The new local ingredient here is the explicit potential at nonempty
bubbling cycles, including its atomic term and its continuity proof.
The algebraic--analytic compactification comparison is supplied by
\cite{BTT17,GSTW21}; we explain the fixed-polarization-class version used
for the Ricci-flat representative.
Once this potential is globalized in the determinant class, uniqueness
in \cite{EGZ09} identifies the metric. The completion theorem of
\cite{Song14} and the orbifold regularity theorem of \cite{GPP25}
then apply. The tangent-cone consequences use
\cite{NO26,DS17,LWX21}; these inputs are cited separately from the local
potential argument.

The normalization of the metric matters throughout. We use
$\|a\|_{L^2}^2=-\int_X\tr(a\wedge *a)$, so a bubble has curvature
energy $8\pi^2$. The rank-normalized Donaldson class is
$-\frac1{2r}p_*\ch_2(\End\mathbb E)$ paired with the surface class.
These conventions determine both $4\pi^2$ in the boundary metric and
$16\pi^4$ in the volume formula.
